\worksheettitle
  {Точка, прямая, отрезок и луч}
  {\modulemark{M01} Версия учителя}

\begin{teacherbox}[Паспорт модуля]
  \textbf{Результат:} ученик различает точку, прямую, отрезок и луч,
  правильно называет их и выполняет построения по обозначениям.

  \textbf{Предварительно:} умеет пользоваться линейкой, различает метку точки
  и её буквенное обозначение. Время не ограничивается; переход определяется
  контрольной точкой.
\end{teacherbox}

\begin{teacherbox}[Смысловой центр]
  Объекты различаются не длиной нарисованной линии, а числом концов и способом
  продолжения. Отдельно проверяется порядок букв в названии луча.
\end{teacherbox}

\section{Вход и разобранный пример}

Во входе: линейка; два конца имеет отрезок. Метки $A$ и $B$ должны быть
видимы отдельно от букв.

\objectscomparisonfigure

\begin{propertybox}[Заполнение примера]
  У прямой нет \textbf{концов}, у отрезка их \textbf{два}, у луча есть
  только \textbf{начало}. Точка обозначается \textbf{заглавной} латинской
  буквой.
\end{propertybox}

\section{Ответы к практике}

\incidenceexamplefigure

\begin{practicebox}[1. Точки и прямая]
  На $a$ лежат $A,B$; не лежит $C$; прямую можно назвать $AB$.
\end{practicebox}

\threelinepointsfigure

\begin{practicebox}[2. Все отрезки]
  $AB,BC,AC$. Если названо только два коротких отрезка, попросите провести
  путь от $A$ до $C$: внутренняя точка $B$ не отменяет отрезок $AC$.
\end{practicebox}

\oppositeraysfigure

\begin{practicebox}[3. Лучи]
  Лучи $AB$ и $AC$, общее начало $A$.
\end{practicebox}

\clearpage
\identifyobjectsanswersfigure

\begin{practicebox}[4. Названия]
  1) точка $M$; 2) прямая $b$; 3) отрезок $CD$; 4) луч $KP$.
\end{practicebox}

\constructionanswersfigure

\begin{teacherbox}[5. Проверка построения]
  Прямая продолжается за обе отмеченные точки. Отрезок заканчивается в $C,D$.
  Луч начинается в $D$, проходит через $A$ и продолжается за $A$.
\end{teacherbox}

\section{Разбор ошибки и контроль}

\begin{warningbox}[Ожидаемое исправление]
  Причина: ученик не учёл роль первой буквы у луча и ошибочно приписал
  отрезку направление. Луч $AB$ начинается в $A$ и проходит через $B$.
  Отрезки $AB$ и $BA$ совпадают, потому что имеют те же два конца.
\end{warningbox}

\begin{checkpointbox}[Ответы]
  \textbf{1.} Луч. \textbf{2.} Неверно: $M,N$ — точки прямой, но не её
  концы. \textbf{3.} Луч $KP$. \textbf{4.} Три отрезка: $DE,EF,DF$.
\end{checkpointbox}

\begin{teacherbox}[Решение о маршруте]
  M01-R: перепутаны число концов, составной отрезок или начало луча.
  M01-H: способ верен, но требуется практика. Устойчивый результат открывает
  M02 и M05; M01\(\star\) необязателен.
\end{teacherbox}
